\documentclass[catalan, a4paper, 12pt, titlepage]{article}
\usepackage{babel}
\usepackage{graphicx}
\usepackage{mathptmx}
%\usepackage{isolatin1}

\title{Una programació didàctica de \\
gestió de bases de dades\\
        \vspace{10mm}
	\includegraphics[width=14cm]{database.eps}
	}

\author{
	Jaume Barceló Vicens\\
	DNI 43135949R\\
	Professors d'ensenyament secudari (0590)\\
	Especialitat informàtica}

\date{
	Família professional informàtica \\
	Cicle formatiu de grau superior d’Administració de Sistemes Informàtics en Xarxa\\
	Mòdul professional de gestió de bases de dades\\
	Codi 0372 170h (+90h d'anglès)}

\begin{document}

\maketitle

\tableofcontents 

\section{Introducció}

Aquesta és la programació didàctica del mòdul ``Gestió de bases de dades'' (codi 0372) del cicle fomatiu de ``Tècnic superior en administració de sistemes i xarxes'' que s'imparteix al CIFP Francesc de Borja Moll el curs escolar 2020-2021.

Aquesta programació tracta de reflexar l'experiència de l'impartició d'aquest mateix mòdul durant el curs 2019-2020. A més, aquesta programació didàctica es presenta com una guia oberta i flexible, sobre la que es poden fer tants canvis com siguin necessaris per adaptar-se a la realitat del grup d'alumnes i del curs en general.

Aquesta programació s'ha desenvolupat tenint en compte la següent normativa estatal:
\begin{itemize}
	\item Llei Orgànica 2/2006, de 3 de maig, d'Educació, que assenyala que el Gobierno de España, prèvia consulta a les comunitats autònomes, establirà les titulacions de formació professional i els aspectes bàsics del currículum.
	\item Llei Orgànica, 8/2013, de 9 de desembre per a la millora de la qualitat educativa i que modifica la LOE.
	\item Llei Orgànica 5/2002, de 19 de juny, de les Qualificacions i de la Formació Professional, que posa en marxa el Sistema Nacional de Qualificacions i Formació professional, desenvolupada pel Reial Decret 1128/2003, de 5 de setemebre, modificat pel Reial Decret 1416/2005, de 25 de novembre, sobre el Catàleg Nacional de Qualificacions Professionals.
	\item Reial Decret 1147/2011, de 29 de juliol, pel qual s'estableix l'ordenació general de la formació professional basada en el Catàleg Nacional de Qualificacions Professionals.
	\item Reial Decret 1629/2009, de 30 d'octubre, pel que s'estableix el títol de ``Tècnic Superior en Administració de Sistemes Informàtics en Xarxa'' i es fixen els seus ensenyaments mínims.
	\item Ordre EDU / 392/2010, de 20 de gener, per la qual s'estableix el currículum del cicle formatiu de grau superior corresponent al títol de ``Tècnic Superior en Administració de Sistemes Informàtics en Xarxa''.
\end{itemize}

En l'àmbit autonòmic considerem la següent normativa:
\begin{itemize}
	\item Decret 91/2012, de 23 de novembre, pel qual s'estableix l'ordenació general de la formació professional del sistema educatiu en el sistema integrat de formació professional a les Illes Balears.
	\item Ordre de la consellera d'Educació i Cultura de 13 de juliol de 2009 per la qual es regula l'organització i el funcionament dels cicles formatius de formació professional del sistema educatiu que s'imparteixen d'acord amb la Llei orgànica 2/2006, de 3 de maig, d'educació, a les Illes Balears, en la modalitat d'ensenyament presencial.
	\item Resolució del conseller d'Educació i Universitat de 18 d'abril de 2018 per la qual s'estableix el calendari escolar del curs 2018-2019 per als centres docents no universitaris de la comunitat autònoma de les Illes Balears.
\end{itemize}

Finalment, en l'àmbit del centre, considerem la següent normativa:
\begin{itemize}
	\item Projecte educatiu de centre (PEC) que planteja les intencions educatives i organitzatives del CIFP: las seva missió, visió i valors. I que inclou com a annexes el projecte lingüístic, el reglament d'organització i funcionament i la concreció curricular del centre.
\end{itemize}

\section{Contextualització}

Aquesta programació didàctica és per al mòdul Gestió de Bases de Dades del Cicle Formatiu de Grau Superior d'Administració de Sistemes Informàtics i Xarxes del Centre Integrat de Formació Professional Francesc de Borja Moll. En aquest centre s'imparteixen també altres cicles formatius de la família d'informàtica i també d'altres famílies professionals.

El centre

\section{Objectius generals}
TBD.

\section{Competències professionals, personals i socials}
TBD.

\section{Resultats d'aprenentatge i criteris d'avaluació}
TBD.

\section{Continguts}
TBD.

\section{Metodologia}
TBD.

\section{Distribució temporal}
TBD.

\section{Avaluació}
TBD.

\section{Recursos materials}
TBD.

\section{Mesures per a l'atenció a la diversitat}
TBD.

\section{Connexió amb altres mòduls}
TBD.

\section{Elements transversals}
TBD.

\section{Unitats didàctiques}
TBD.

  \subsection{UD 1}
  TBD.

  \subsection{UD 2}
  TBD.

  \subsection{UD 3}
  TBD.

  \subsection{UD 4}
  TBD.

  \subsection{UD 5}
  TBD.

  \subsection{UD 6}
  TBD.

  \subsection{UD 7}
  TBD.

  \subsection{UD 8}
  TBD.

  \subsection{UD 9}
  TBD.

  \subsection{UD 10}
  TBD.

  \subsection{UD 11}
  TBD.

  \subsection{UD 12}
  TBD.

\section{Conclusió}
TBD.

\section{Bibliografia}
TBD.

...

Això és una prova.
\end{document}
